\chapter{Questions, negotiation, retrospective}
\label{ch:after}

\section*{What this chapter is for}

Three things that happen at the edges of a process and are usually improvised:
the questions you ask, the conversation about the offer, and what you keep
afterwards. The third is the one that compounds.

\section{The questions you ask}

Your questions are assessed, and they are assessed for what they reveal about
what you consider important.

\begin{kittable}{@{}XX@{}}
\textbf{Reveals little} & \textbf{Reveals something} \\
\midrule
``What does a typical day look like?'' & ``What did the last person in this role spend most of their time on, and was that the plan?'' \\
``What is the tech stack?'' & ``What in the stack are you least happy with, and what would it take to change it?'' \\
``How do you handle work-life balance?'' & ``When was the last time this team shipped something late, and what happened?'' \\
``What are the growth opportunities?'' & ``What does somebody have to do here to move from this level to the next, and who has done it recently?'' \\
\end{kittable}

The right column works because each question has a wrong answer, which means
the answer carries information. \judgement{A question that cannot be answered
badly is not a question, it is a pleasantry, and asking three of them in a row
is itself a signal.}

Two worth asking at this level specifically, both of which come straight from
earlier chapters: what the hardest currently-unsolved problem in the team is
(Chapter~\ref{ch:person} \dash{} it gives you the calibration you were looking
for), and how the team knows when one of its model-driven features has got
worse (Chapter~\ref{ch:evals} \dash{} the answer tells you whether they have
been hurt yet, and it is the most informative question in this domain).

\section{Negotiation, briefly}

This book carries no figures and no scripts. Three structural points, which are
what generalise.

\textbf{Where a company runs a tiered ladder, the level is the negotiation.}
The band follows from the row, so a conversation about the number is usually a
conversation about the row, and it is far easier to have before the offer than
after it.

\textbf{Establish composition before discussing a figure.} Base, variable,
equity and their vesting are four different things and a single headline number
conceals which mix is on offer. Chapter~\ref{ch:company} puts this in the
screen, where it belongs.

\textbf{What is movable varies by company, not by candidate.} Start date and
scope are often flexible; the band ceiling for a row usually is not. Asking
early which is which costs nothing and prevents the late-stage version of the
conversation, which is the expensive one.

\section{The retrospective}

Write one after every process, including the ones that end early, and
especially the ones that end early.

\begin{kitmethod}
\textbf{What happened} \dash{} stages, dates, who, outcome. Facts only.

\textbf{What the process measured} that you had not expected. This is where
most of the value is: a stage that turned out to be about something else is
information about the market, not about that company.

\textbf{What you got wrong in your estimate,} with the numbers you wrote down
at the time. If you did not write them down, that is the first finding.

\textbf{Which class it ended in} \dash{} Chapter~\ref{ch:classes}. Above the
line or below it. This single field, accumulated over several processes, is
the most useful thing in the whole document.

\textbf{What you would reuse,} and what you would not build again.
\end{kitmethod}

The structure is \code{templates/retrospective.md}. The running table across
processes is \code{templates/outcome-ledger.md}: one row each, with the stage
reached and the class it ended in.

\begin{kitnote}
The ledger is the point, not the individual retrospectives. One process tells
you almost nothing \dash{} outcomes at this level are noisy and a single
rejection is weak evidence about anything. Four processes with a pattern in the
last column tell you where your actual constraint is, and that is not knowable
any other way.

If three of four ended below the line, no amount of further preparation is the
answer, and the ledger is what makes that visible rather than arguable.
\end{kitnote}

\section{After a rejection}

Short, specific, no request to reconsider. Thank them for being direct if they
were, since most are not. If the reason given was a data objection, saying that
you understand it is more credible than arguing, and it costs nothing.

\judgement{Leave a door open only if you would actually walk through it. A
sentence inviting future contact about a specific kind of work is worth
sending; a general offer to be considered for anything is not, and reads as
what it is.}

\begin{drill}
Write the retrospective for your most recent process, even if it was months
ago, and even if it went well. Fill in the class field honestly. Then start the
ledger with that one row \dash{} the document only becomes useful at three or
four rows, and it never gets started at the point when it would be useful.
\end{drill}

\section*{What this chapter does not claim}

It does not claim particular questions produce particular outcomes. The claim
is about information: a question with a possible bad answer tells you something
and a question without one does not.

No negotiation figures, tactics or scripts are offered. The three structural
points are about how tiered compensation systems are built, and they do not
apply to companies that do not use one.
